\documentclass[12pt]{article}
\usepackage{amsmath}
\usepackage{amssymb}
\title{From Newton's Laws to Orbital Mechanics}
\author {Sergey}
\date{\today}

\begin{document}
\maketitle

\begin{abstract}
This paper derives the fundamental results of orbital mechanics
from Newton's three laws of motion and his law of universal gravitation. To many, this will already be obvious intro to physics, but to me, this is entirely new and I had
to learn all of this from scratch. This entire "paper" is just me teaching myself basic physics, how to use LaTeX, navigating in neovim, and generally having fun. 
I plan to make this, amongst other similar "papers", a series I call "Documented Learning." Over time, I hope to document the different things I learned over the course 
of my self-study with the hope that I can look back one day and say "Wow. I really didn't even know Newton's laws of motion at some point, huh."
\end{abstract}

\newpage
\section{Newton's Laws of Motion}
First and foremost obvious is that an object at rest will stay at rest, and an object in motion will remain in motion unless acted upon by some external force. 

More interesting is that the total force on an object is the product of its mass and acceleration. $F = ma$. Speed is the measure of how fast an object is moving, 
whereas acceleration is the measure of how fast the speed of said object is changing.
Basic algebra will teach that if $F$ is constant, increasing either $m$ or $a$ will lead to a decrease in the other.

For the third law, every action has an equal and opposite reaction. The Earth pulls the moon with the same strength that the moon pulls the Earth. 

\section{Universal Gravitation}
This measures how strong gravity pulls at one specific distance, measured in Newtons. Every object with a mass attracts every other object with a mass 
per the formula: $$F = G \cdot \frac{M_1 M_2}{D^2}$$ where G is the gravitational constant, M1 and M2 are the masses of each object, and D is the distance 
between the centerpoints of both masses. 

\section{Gravitational Acceleration}
This measures how much acceleration the gravitational force of the primary body produces on the secondary body. Given $F = ma$, I can actually set this equal to 
$F = G \frac{M_1 M_2}{D^2}$, where: $$ma = G \cdot \frac{M_1 M_2}{D^2}$$ 
In this formula, I'm ultimately solving for $a$, which would lead both sides to be divided by $m$ in: $$a = G \cdot \frac{M_1 M_2}{D^2 m}$$ 
Where $m$ substitutes the mass being pulled. I actually wondered why it mattered - that is, whether $m$ should cancel $M_1$ or $M_2$, 
but it doesn't. $M$ is the primary body, and $m$ is the secondary body. Basic algebra teaches that the $m$'s cancel out, thus leaving the $a$ part of this formula, 
or gravitational acceleration, to be: $$ a = G \cdot \frac{M}{D^2}$$ When $m$ cancels, acceleration is entirely dependent only upon $M$ and $D$. A bowling ball and a 
feather, for example, fall at the same rate towards the Earth, where Earth constitutes $M$ and either ball or feather constitutes $m$.

\section{Escape Velocity}
The way I interpret escape velocity is the speed needed to escape the gravitational force of a primary body. Before I can derive some given escape velocity, I 
learned a bit about how the formula for kinetic energy was discovered. Funnily enough, it was basically discovered experimentally by Leibniz, and is:
$$\frac{1}{2}M_2v^2$$
I then learned a bit about the interaction between kinetic and potential energy, and for any object in motion, they should basically always be constant,
which is the law of conservation of energy. In the previous equation, $M_2$ is mass of said object, $v$ is velocity, which is what I am calculating for in escape velocity, and
the 0.5 is some sort of energy integral that I am too inexperienced to derive yet. I'll take it as a given and come back to it in a future project or white paper.

For potential energy, the formula is:
$$- G \cdot \frac{M_1 M_2}{D}$$
Negative because the reference point is basically the "top" of the primary body's gravity field. $M_1$ remains the primary body, $M_2$ the secondary body, and $D$ the distance
between them. On Earth, potential energy is always negative, kinetic energy is always positive. The sum of the two is total energy. The way I understand this is that if 
KE exceeds PE for an object, that object can escape the gravitational bound of the primary body. If PE exceeds KE, then it cannot ever escape the gravitational
bound of said primary body. Interestingly, KE can equal PE, and that's escape velocity. Of course, that would be absolute value of PE given it's always negative.
I'm not sure why but I struggled a bit understanding this concept, even though it's incredibly easy in retrospect. 
Regardless, velocity $v$ can be solved from the total energy formula, which is: $$\frac{1}{2}M_2v^2 + - G \cdot \frac{M_1 M_2}{D} = 0$$
To be clear, the $= 0$ part of the above equation is because KE is equal to the absolute value of PE, which means added they must be 0. And, that formula rearranged, is:
$$v_{escape} = \sqrt{2G \cdot \frac{M_1}{D}}$$

\section{Orbital Velocity}
Orbital velocity as I understand it is just the speed at which a secondary body orbits a primary body. For orbital velocity, 
I needed to learn about centripetal force. Anything force I already know is $F = ma$, so I can infer centripetal force is related somehow
to centripetal acceleration. I learned about the proper definition of acceleration, which is that it's a change in direction AND speed (velocity), so I imagine some sort 
of satellite orbiting the Earth at a constant speed but constantly changing direction due to gravity pulling it to the center of Earth. That would still be 
considered acceleration. The reason the speed of the satellite never changes is because it was launched at some velocity and it just keeps going in said velocity, which 
is Newton's first law. I learned that this acceleration, which is centripetal acceleration in this case, can only increase or decrease as a result of two things: 
either the satellite moves incredibly fast, meaning it has a high velocity $v$, and changes direction quickly as well. I would imagine this is a 
satellite that performs many revolutions around Earth. But it could also decrease with distance $D$, where a satellite on the very edge of the orbit of Earth must,
by definition, perform revolutions around Earth slower than a satellite closer to the Earth assuming they are moving at the same velocity. So, the formula for centripetal 
acceleration should be some proportion of $\frac{M_2v}{D}$, and is: $$\frac{M_2v^2}{D}$$
I learned that understanding the $v^2$ also includes calculus I don't quite understand yet, so I'll come back to that later and accept the formula as a given for now.
With that given, in any given orbital velocity, we can assume the centripetal force comes from gravitational force. That only makes sense, and I think this is 
definitional, but I can't prove it. Regardless, we can take that as a given and set
$$G \cdot \frac{M_1 M_2}{D^2} = \frac{M_2v^2}{D}$$ 
Where $M_2$ is the orbiting body. And after basic algebra, and solving for orbital velocity $v$, gives:
$$v_{orbital} = \sqrt{G \cdot \frac{M_1}{D}}$$
To clarify, since the mass of the secondary body (the orbiting body) gets cancelled out in the formula, its mass is actually irrelevant in this calculation.

\end{document}